\documentclass[11pt]{article}

\usepackage[margin=1in]{geometry}
\usepackage{amsmath}
\usepackage{booktabs}
\usepackage{enumitem}
\usepackage{graphicx}
\usepackage{siunitx}
\usepackage{hyperref}
\usepackage{gensymb}


\title{Basic Static Test}
\author{Rift / Team 09}
\date{\today}

\begin{document}


\maketitle

\section*{Test ID}
\textbf{T-01BS Basic Static Test} 

\section*{Test Objective}
Qualify the capability of the rover to transition from prone to standing position. Prone is defined as the configuration with the payload in contact with the ground and the truss legs in equilateral triangle position. Standing is determined to be the configuration with the payload not in contact with the ground and the necessary truss shape to hold the payload in this static, elevated position. Quantify the maximum weight that the payload can handle before truss leg failure.

\section*{Robot Configuration}
\begin{itemize}[noitemsep]
    \item Robot configurations: Prone and Standing
    \item Tube pressures: 15 psi
    \item Control mode: Open-loop
    \item Electronics status: Batteries are fully charged
\end{itemize}

\section*{Materials}
\begin{itemize}[noitemsep]
    \item Plywood payload container (mass = ...)
    \item Weights (incrementing by 1 kg, 5 kg, 10 kg)
    
\end{itemize}

\section*{Environment and Setup}
\begin{itemize}[noitemsep]
    \item Surface type: Finished Smooth Concrete
    \item Surface angle: $\psi = 0 {\degree}$ 
    \item External load: None 
    \item Measurement method: increasing weights and recording capacity
\end{itemize}

\section*{Initial Conditions}
\begin{itemize}[noitemsep]
    \item Initial position: $x_0 = 0$
    \item Initial orientation: $\theta_z = 0 {\degree}$
    \item Initial shape state: Prone
    \item Stabilization time before command: 10?~s
\end{itemize}

\section*{Command Sequence}
\begin{itemize}[noitemsep]
    \item Command type: \underline{\hspace{5 cm}}
    \item Command values: \underline{\hspace{5 cm}}
    \item Update rate: 3~Hz
    \item Command Steps: 8~s
\end{itemize}

\section*{Measured Quantities}
\begin{itemize}[noitemsep]
    \item Total weight supported by payload
\end{itemize}

\section*{Success Criteria (Per Trial)}
\begin{itemize}[noitemsep]
    \item Total weight supported: 35 kg  $\leq$ 45 kg $\leq$ No upper limit 
    \item Tube pressure $\geq$ 12 psi
    \item The payload does not return to the prone position under load
\end{itemize}

\section*{Test Procedure}
\begin{enumerate}[noitemsep]
    \item Place robot in prone position
    \item Complete pretest checklist
    \item Send command to move rover from prone to standing position
    \item Place masses, increasing from 1 kg increments to 5 kg to 10 kg
    \item Record maximum mass that the rover can support before static failure (buckling)

\end{enumerate}

\section*{Number of Trials}
\begin{itemize}[noitemsep]
    \item Total trials: 1
    \item Time between trials: N/A
\end{itemize}

\section*{Results}
Summarize the observed performance metrics.

\begin{center}
\begin{tabular}{lc}
\toprule
Metric & Value\\
\midrule
Mass (kg)&  \\
\end{tabular}
\end{center}

\section*{Observed Failure Modes}
Describe any observed failure behaviors.

\begin{itemize}[noitemsep]
    \item Truss buckles at ... kg
    \item The payload lower edge displaces ... cm from standing position
\end{itemize}

\section*{Conclusions}
Provide a concise interpretation of the results.

\textit{Example:} The robot fully lifted itself from a prone position. The robot was loaded to 35 kg before a tube buckled.

\section*{Notes and Limitations}
Document sources of variability or limitations.

\end{document}